% =====================================================================
%  CHAPITRE 17 — LECTURE ET ÉCRITURE
% =====================================================================
\chapter{Lire et écrire : orthographe, ponctuation, modèles}

\section{Partie A --- La lecture}

\subsection{Pourquoi lire est l'accélérateur le moins fatigant}
Lire apporte le vocabulaire en contexte, la grammaire par imprégnation, et la structure des
phrases anglaises naturelles. Vingt minutes de lecture agréable par jour valent une heure
d'exercices pénibles. Mais il faut lire \textbf{au bon niveau} --- c'est là que tout se joue.

\begin{regle}[La règle des 98\,\%]
Prenez une page au hasard. Comptez les mots que vous ne connaissez pas :\\
\textbf{0 à 2 mots inconnus sur 100} \ra niveau idéal, lisez pour le plaisir.\\
\textbf{3 à 5 mots} \ra niveau utile mais fatigant : travaillez 10 minutes, pas plus.\\
\textbf{Plus de 5 mots} \ra trop difficile : changez de livre, ce n'est pas une question
de volonté.
\end{regle}

\begin{tableau}[Les deux lectures, à alterner]
\begin{tabularx}{\linewidth}{@{}p{3.2cm}p{5.2cm}X@{}}
\entete{} & \entete{Lecture extensive} & \entete{Lecture intensive} \tabularnewline
\hline
Volume & 20 à 40 pages faciles par jour & 1 à 2 pages difficiles par jour \tabularnewline
But & Plaisir, fluidité, volume de mots rencontrés & Précision, vocabulaire actif, structures \tabularnewline
Dictionnaire & \textbf{Interdit} (ou une fois par page) & Autorisé, avec notes \tabularnewline
À voix haute & Non & Oui, en imitant le rythme \tabularnewline
Ce qu'on retient & Le sens général & Chaque mot, chaque tournure \tabularnewline
\end{tabularx}
\end{tableau}

\subsection{Que lire, mois par mois}
\begin{tableau}[Parcours de lecture sur 12 mois]
\begin{tabularx}{\linewidth}{@{}p{1.8cm}p{4.6cm}X@{}}
\entete{Mois} & \entete{Matériau} & \entete{Conseil} \tabularnewline
\hline
1-2 & Textes de ce livre ; dialogues ; listes de phrases ; panneaux, étiquettes, notices d'appareils & Relisez ce que vous connaissez déjà : la relecture crée la vitesse. \tabularnewline
3-4 & Livres simplifiés (\en{graded readers}) niveau A2 ; livres pour enfants illustrés ; articles courts en ligne & 10 pages par jour, sans dictionnaire. Un livre entier en une semaine. \tabularnewline
5-6 & \en{Graded readers} niveau B1 ; nouvelles courtes ; articles de presse simples ; bandes dessinées & Notez 10 expressions par semaine, pas 50 mots. \tabularnewline
7-8 & Premiers romans faciles (jeunesse, policier simple) ; articles d'actualité ; blogs sur vos centres d'intérêt & Changez de genre chaque semaine. \tabularnewline
9-10 & Romans pour adultes, presse (actualités, sport, économie), essais courts & Lisez le même sujet dans deux sources différentes. \tabularnewline
11-12 & Littérature, presse d'opinion, documentation professionnelle de votre métier & Objectif : lire comme vous lisez en français --- vite, et pour le contenu. \tabularnewline
\end{tabularx}
\end{tableau}

\begin{astuce}[Choisissez des livres que vous avez déjà lus en français]
Relire en anglais un roman ou une bande dessinée que vous connaissez par c\oe ur est la
méthode la plus efficace qui existe : vous comprenez l'histoire, donc votre cerveau se
concentre sur la \textbf{forme}. Même chose avec les films : revoir un film déjà vu avec
sous-titres anglais fait des miracles.
\end{astuce}

\subsection{Trois textes de lecture commentés}

\subsubsection*{Texte A --- niveau A2 : une lettre d'un ami}
\begin{dialoguebox}[Texte A]
\begin{anglais}
Dear Sam,

Thanks for your message --- it was great to hear from you. Sorry I haven't written for so
long, but work has been really busy these last few weeks.

Let me tell you my news. In March, I moved to a new flat in the centre of town. It's small,
but it has a balcony, and I can walk to work in fifteen minutes. I've been living there for
two months now and I'm very happy with it.

I've also started playing football again, twice a week with some colleagues. I'm not as fit
as I used to be, so the first games were hard! Last Saturday we won three-two, and I scored
--- the first goal of my life, believe me.

What about you? Are you still working at the hospital? How are your parents? I'm planning to
visit Dakar in August, so maybe we could meet up and have dinner.

Write soon and send me some photos.

Best wishes,\\
Moussa
\end{anglais}
\end{dialoguebox}

\textbf{Questions :} 1. Pourquoi Moussa n'a-t-il pas écrit plus tôt ? 2. Depuis combien de
temps habite-t-il son nouvel appartement ? 3. Quel avantage principal présente cet
appartement ? 4. Pourquoi les premiers matchs ont-ils été difficiles ?
5. Que propose-t-il pour août ? 6. Relevez trois temps différents et expliquez leur usage.

\subsubsection*{Texte B --- niveau B1 : un article}
\begin{dialoguebox}[Texte B]
\begin{anglais}
\textbf{Why walking is the most underrated exercise}

Most people think that to get fit, they need a gym, expensive equipment and several hours a
week. In fact, one of the best things you can do for your health costs nothing and requires
no training: walking.

Studies have shown that thirty minutes of brisk walking a day reduces the risk of heart
disease, improves sleep and lowers stress. Unlike running, walking can be continued into old
age, and it rarely causes injuries. It also has an advantage that is often forgotten: you
can do it with other people, which makes it easier to keep going.

So how do you start? Forget about the idea of a daily hour-long walk --- it will last two
weeks. Instead, add walking to what you already do. Get off the bus one stop earlier. Walk
while you are on the phone. Take the stairs instead of the lift. After a month, you will
have added forty minutes a day to your life without ever finding the time.

The hardest part of exercise is not the exercise. It is starting, and continuing.
\end{anglais}
\end{dialoguebox}

\textbf{Questions :} 1. Quelle idée reçue l'article conteste-t-il ? 2. Citez trois bénéfices
de la marche. 3. Pourquoi la marche est-elle préférable à la course selon l'auteur ?
4. Quel conseil concret donne-t-il pour commencer ? 5. Que signifie la dernière phrase ?
6. Repérez deux exemples de voix passive et un comparatif.

\subsubsection*{Texte C --- niveau B2 : un extrait argumentatif}
\begin{dialoguebox}[Texte C]
\begin{anglais}
\textbf{Should English replace French as our second language?}

It is a question that comes up every few years, usually after a diplomatic meeting or a
football match, and it deserves a calmer answer than the one it usually receives.

The argument for English is, on the surface, irresistible. It is the language of science, of
aviation, of computing, of most international business and of a very large share of the
entertainment the world consumes. A young person who speaks English well has access to
information that a young person who does not simply cannot reach.

But the argument is often presented as though languages were tools in a box, to be kept or
discarded according to usefulness. They are not. A language is also a way of thinking, a
literature, a set of relationships, and in our region, a history. French connects us to more
than twenty countries and to a diplomatic tradition in which we have invested for decades.
Dropping it would not make us more competitive; it would make us poorer.

The sensible answer is therefore not either/or, but both --- and, increasingly, more. The
real problem is not that we speak too many languages, but that we speak them badly. Improving
English teaching should not mean abandoning French; it should mean being honest about how
little of either we currently master.
\end{anglais}
\end{dialoguebox}

\textbf{Questions :} 1. Quand cette question resurgit-elle habituellement ? 2. Quels sont les
domaines où l'anglais domine ? 3. Quelle métaphore l'auteur rejette-t-il, et pourquoi ?
4. Que perdrait-on en abandonnant le français ? 5. Quelle est la conclusion ?
6. Que critique l'auteur dans la dernière phrase ? 7. Relevez trois connecteurs logiques et
expliquez leur fonction.

\begin{corrige}[Textes A, B, C]
\textbf{A.} 1. Le travail a été très chargé. 2. Deux mois. 3. Il est à quinze minutes à pied
du travail (et il a un balcon). 4. Il n'est plus aussi en forme qu'avant (\en{not as fit as
I used to be}). 5. Se rencontrer et dîner ensemble à Dakar. 6. \en{I haven't written}
(\en{present perfect} : lien avec le présent), \en{I moved} (passé simple : date précise),
\en{I've been living there for two months} (\en{present perfect continuous} : durée en cours),
\en{I used to be} (habitude passée révolue).\\[4pt]
\textbf{B.} 1. L'idée qu'il faut une salle de sport, du matériel coûteux et beaucoup de temps.
2. Réduit le risque de maladie cardiaque, améliore le sommeil, diminue le stress.
3. On peut la continuer en vieillissant et elle cause rarement des blessures ; on peut aussi
la pratiquer à plusieurs. 4. Ajouter la marche à ce qu'on fait déjà : descendre un arrêt plus
tôt, marcher au téléphone, prendre les escaliers. 5. Le plus dur n'est pas l'effort, c'est de
commencer et de continuer. 6. Passifs : \en{that is often forgotten}, \en{it will be
continued} (ou \en{can be continued}) ; comparatif : \en{Unlike running\ldots} /
\en{easier to keep going}.\\[4pt]
\textbf{C.} 1. Après une réunion diplomatique ou un match de football. 2. Science, aviation,
informatique, affaires internationales, divertissement. 3. La métaphore des langues comme
outils dans une boîte : une langue est aussi une pensée, une littérature, des relations, une
histoire. 4. Des relations avec plus de vingt pays et une tradition diplomatique ; on
deviendrait plus pauvre, pas plus compétitif. 5. Non pas l'un ou l'autre, mais les deux ---
et même davantage. 6. Nous parlons mal les deux langues, plutôt que trop de langues.
7. \en{therefore} (conséquence), \en{but} (opposition), \en{so} (conséquence),
\en{unlike} (contraste), \en{increasingly} (progression).
\end{corrige}


\section{Partie B --- L'orthographe}

\subsection{Les sept règles qui résolvent 90\,\% des problèmes}
\begin{tableau}[Règles d'orthographe]
\begin{tabularx}{\linewidth}{@{}p{4.2cm}X@{}}
\entete{Règle} & \entete{Application} \tabularnewline
\hline
\textbf{1. Consonne doublée} & Verbe d'une syllabe, voyelle + consonne : on double avant
\en{-ing} et \en{-ed} : \en{stop} \ra \en{stopped, stopping} ; \en{plan} \ra \en{planned,
planning} ; \en{travel} \ra \en{travelled} (BrE) / \en{traveled} (AmE). \tabularnewline
\textbf{2. \en{e} final supprimé} & Devant une voyelle : \en{make} \ra \en{making} ;
\en{write} \ra \en{writing} ; mais conservé devant consonne : \en{making} / \en{likes}. \tabularnewline
\textbf{3. \en{y} \ra \en{i}} & Après consonne : \en{study} \ra \en{studies, studied} ;
\en{happy} \ra \en{happier, happiness}. Après voyelle, on garde : \en{play} \ra \en{plays,
playing}. \tabularnewline
\textbf{4. \en{i} avant \en{e}} & Sauf après \en{c} : \en{believe, receive, piece, ceiling}.
Exceptions : \en{weird, height, either, their}. \tabularnewline
\textbf{5. Pluriels en \en{-es}} & Après \en{-s, -sh, -ch, -x, -o} : \en{boxes, brushes,
watches, buses, tomatoes}. \tabularnewline
\textbf{6. \en{-f} \ra \en{-ves}} & \en{leaf} \ra \en{leaves}, \en{life} \ra \en{lives},
\en{wife} \ra \en{wives}, \en{knife} \ra \en{knives} ; mais \en{roofs, chiefs, beliefs}. \tabularnewline
\textbf{7. Adverbes en \en{-ly}} & \en{happy} \ra \en{happily} ; \en{true} \ra \en{truly} ;
\en{terrible} \ra \en{terribly} ; mais \en{public} \ra \en{publicly}. \tabularnewline
\end{tabularx}
\end{tableau}

\subsection{Orthographe britannique ou américaine ?}
\begin{memo}[Choisissez-en une et gardez-la]
Mélanger les deux dans un même document est la seule vraie faute. En Afrique de l'Ouest,
l'enseignement suit généralement le britannique ; les médias et Internet sont majoritairement
américains. Recommandation : \textbf{britannique à l'écrit formel} (école, administration,
examens), \textbf{américain accepté} dans les échanges informels et sur Internet.
\end{memo}

\begin{tableau}[Les 40 différences les plus courantes]
\footnotesize
\begin{longtable}{@{}p{3.4cm}p{3.4cm}p{3.4cm}p{3.4cm}@{}}
\entete{Britannique} & \entete{Américain} & \entete{Britannique} & \entete{Américain} \tabularnewline
\endfirsthead
\entete{BrE} & \entete{AmE} & \entete{BrE} & \entete{AmE} \tabularnewline
\endhead
\hline
\en{colour} & \en{color} & \en{organise} & \en{organize} \tabularnewline
\en{favourite} & \en{favorite} & \en{realise} & \en{realize} \tabularnewline
\en{honour} & \en{honor} & \en{apologise} & \en{apologize} \tabularnewline
\en{behaviour} & \en{behavior} & \en{analyse} & \en{analyze} \tabularnewline
\en{neighbour} & \en{neighbor} & \en{practise} (v.) & \en{practice} (v.) \tabularnewline
\en{centre} & \en{center} & \en{licence} (n.) & \en{license} (n.) \tabularnewline
\en{theatre} & \en{theater} & \en{defence} & \en{defense} \tabularnewline
\en{metre} & \en{meter} & \en{offence} & \en{offense} \tabularnewline
\en{litre} & \en{liter} & \en{programme} & \en{program} \tabularnewline
\en{travelled} & \en{traveled} & \en{catalogue} & \en{catalog} \tabularnewline
\en{cancelled} & \en{canceled} & \en{dialogue} & \en{dialog} \tabularnewline
\en{labelled} & \en{labeled} & \en{cheque} & \en{check} (bancaire) \tabularnewline
\en{fulfil} & \en{fulfill} & \en{grey} & \en{gray} \tabularnewline
\en{enrol} & \en{enroll} & \en{tyre} & \en{tire} (pneu) \tabularnewline
\en{jewellery} & \en{jewelry} & \en{aeroplane} & \en{airplane} \tabularnewline
\en{kerb} & \en{curb} & \en{lorry} & \en{truck} \tabularnewline
\en{pavement} & \en{sidewalk} & \en{petrol} & \en{gas / gasoline} \tabularnewline
\en{lift} & \en{elevator} & \en{flat} & \en{apartment} \tabularnewline
\en{biscuit} & \en{cookie} & \en{sweets} & \en{candy} \tabularnewline
\en{holiday} & \en{vacation} & \en{autumn} & \en{fall} \tabularnewline
\end{longtable}
\end{tableau}

\section{Partie C --- La ponctuation et la typographie anglaises}
\begin{tableau}[Douze règles de ponctuation --- différentes du français]
\begin{tabularx}{\linewidth}{@{}p{4.6cm}X@{}}
\entete{Règle} & \entete{Détail} \tabularnewline
\hline
\textbf{Pas d'espace avant} \en{, ; : ! ?} & Contrairement au français : \en{Hello! How are you?}
(aucune espace avant \en{!} ni \en{?}). \tabularnewline
\textbf{Guillemets} & \en{``double quotes''} en anglais américain et souvent britannique ;
\en{`single quotes'} pour une citation dans la citation. Pas de « ». \tabularnewline
\textbf{Apostrophe} & \en{'} pour la possession (\en{Ali's book}) et les contractions
(\en{don't, I'm}). \textbf{Jamais} de pluriel avec apostrophe : \en{the 1990s} (pas \en{1990's}). \tabularnewline
\textbf{Majuscules} & Jours, mois, langues, nationalités prennent une majuscule :
\en{Monday, January, English, Senegalese}. Pas les adjectifs de matière ni les titres
complets (\en{the president} sauf devant un nom : \en{President Diop}). \tabularnewline
\textbf{Virgule} & Avant \en{and} dans une énumération (virgule d'Oxford) :
\en{bread, milk, and eggs}. Après un connecteur en début de phrase : \en{However, I disagree.} \tabularnewline
\textbf{Point-virgule} & Plus rare qu'en français ; préférer deux phrases. \tabularnewline
\textbf{Deux-points} & Pour introduire une liste ou une explication : \en{I need three things:
time, patience and coffee.} \tabularnewline
\textbf{Tiret long} & \en{---} (em dash) pour une incise, sans espaces en américain :
\en{He arrived --- late, as usual --- and apologised.} \tabularnewline
\textbf{Nombres} & Virgule pour les milliers, point pour les décimales :
\en{12,500 francs} ; \en{3.5 kilometres} --- \textbf{l'inverse du français}. \tabularnewline
\textbf{Heures} & \en{6.30 pm} (BrE) ou \en{6:30 pm} (AmE) ; \en{a.m. / p.m.} \tabularnewline
\textbf{Dates} & \en{5 May 2026} (BrE) ou \en{May 5, 2026} (AmE). \tabularnewline
\textbf{Titres} & Majuscule aux mots importants : \en{How to Learn English in Twelve Months}. \tabularnewline
\end{tabularx}
\end{tableau}

\begin{piege}[Les cinq erreurs de ponctuation les plus visibles]
\en{*Hello !} \ra \en{Hello!} \qquad
\en{*« I'm fine », he said.} \ra \en{``I'm fine,'' he said.} \qquad
\en{*12.500 francs} \ra \en{12,500 francs} \qquad
\en{*the 1990's} \ra \en{the 1990s} \qquad
\en{*in monday} \ra \en{on Monday}
\end{piege}

\section{Partie D --- Écrire : huit modèles prêts à l'emploi}

\subsection{Modèle 1 --- E-mail informel}
\begin{dialoguebox}[E-mail à un ami ou un partenaire de langue]
\begin{anglais}
\textbf{Subject: Let's practise this week?}

Hi John,

Hope you're doing well! Thanks for your last message --- it really helped.

I've been working on pronunciation this month, and I think I've improved a bit, but I still
struggle with the ``th'' sound. Would you like to do a voice call this weekend? I'm free on
Saturday afternoon or Sunday morning, your time.

In exchange, I can help you with your French --- I've prepared a short list of expressions
that people actually use.

Let me know what works for you.

Take care,\\
Moussa
\end{anglais}
\end{dialoguebox}
\textbf{Structure :} salutation courte \ra une phrase de lien \ra l'objet de l'e-mail
\ra la proposition précise \ra l'offre en retour \ra formule de fin chaleureuse.

\subsection{Modèle 2 --- E-mail professionnel (demande)}
\begin{dialoguebox}[Demande d'information à une entreprise]
\begin{anglais}
\textbf{Subject: Enquiry about the training programme --- ref. 2026/EN}

Dear Mr Mensah,

I am writing to enquire about the digital marketing training programme advertised on your
website last week.

I would be grateful if you could provide me with the following information:
\begin{itemize}
  \item the dates and duration of the next session;
  \item the total cost, and whether payment can be made in instalments;
  \item whether the course is available online for participants based in Thiès.
\end{itemize}

I have attached my CV for your reference. I look forward to hearing from you.

Yours sincerely,\\
Moussa Diop\\
\textit{Accountant --- Tel: +221 77 000 00 00}
\end{anglais}
\end{dialoguebox}
\textbf{Les blocs fixes à retenir :} \en{I am writing to\ldots} /
\en{I would be grateful if you could\ldots} / \en{Please find attached\ldots} /
\en{I look forward to hearing from you} (+ \en{-ing} !) /
\en{Yours sincerely} (si vous connaissez le nom) --- \en{Yours faithfully} (si
\en{Dear Sir or Madam}).

\subsection{Modèle 3 --- Réclamation polie mais ferme}
\begin{dialoguebox}[Produit défectueux]
\begin{anglais}
\textbf{Subject: Faulty product --- order \#45892}

Dear Sir or Madam,

On 12 May I ordered a wireless printer from your online shop (order number 45892). It was
delivered on 18 May, but unfortunately it does not work properly: the printer turns off after
a few minutes and cannot connect to my computer.

I have already contacted your support team twice by phone, but the problem has not been
resolved. I am therefore writing to request either a replacement or a full refund.

I would appreciate a reply within seven days. I have attached a copy of the invoice and a
short video showing the problem.

Yours faithfully,\\
Fatou Ndiaye
\end{anglais}
\end{dialoguebox}
\textbf{La technique du sandwich appliquée à l'écrit :} faits datés \ra problème précis
\ra démarches déjà faites \ra demande claire (\en{either\ldots or\ldots}) \ra délai
\ra pièces jointes \ra formule polie. Aucune émotion, aucun reproche personnel : c'est ce
qui rend une réclamation efficace.

\subsection{Modèle 4 --- Lettre de motivation (structure en 4 paragraphes)}
\begin{dialoguebox}[Squelette à personnaliser]
\begin{anglais}
\textbf{Paragraph 1 --- Why I am writing.}\\
\en{I am writing to apply for the position of \ldots, which I saw advertised on \ldots\ I am
very interested in this role because \ldots}

\textbf{Paragraph 2 --- Why I am suitable (2 proofs + figures).}\\
\en{I have \ldots years' experience in \ldots\ In my current position at \ldots, I am
responsible for \ldots\ Last year, I managed to \ldots, which resulted in a \ldots per cent
increase in \ldots}

\textbf{Paragraph 3 --- Why this company.}\\
\en{I am particularly attracted to your company because of \ldots\ I believe my experience in
\ldots would be a valuable addition to your team.}

\textbf{Paragraph 4 --- Availability and closing.}\\
\en{I would welcome the opportunity to discuss my application in an interview. I am available
from \ldots\ Thank you for considering my application. Yours sincerely, \ldots}
\end{anglais}
\end{dialoguebox}

\subsection{Modèle 5 --- CV anglais (les différences à connaître)}
\begin{tableau}[CV britannique/américain : ce qu'on met et ce qu'on ne met pas]
\begin{tabularx}{\linewidth}{@{}p{5.0cm}X@{}}
\entete{Élément} & \entete{Règle} \tabularnewline
\hline
Photo & \textbf{Pas de photo} dans les pays anglophones (sauf demande explicite). \tabularnewline
Âge, situation familiale, religion & \textbf{À ne jamais mentionner} --- c'est illégal ou très mal vu. \tabularnewline
Titre & \en{Accountant with 6 years' experience in the retail sector} --- un titre avec une valeur, pas seulement un métier. \tabularnewline
Profil (\en{Personal statement}) & 3 lignes en haut du CV : qui vous êtes, ce que vous savez faire, ce que vous cherchez. \tabularnewline
Expérience & Ordre \textbf{anti-chronologique} ; puces commençant par un \textbf{verbe d'action} : \en{Managed, Developed, Increased, Organised, Trained, Reduced}. \tabularnewline
Résultats & Toujours chiffrés : \en{increased sales by 20\,\%}, \en{trained a team of 12}, \en{reduced costs by 3 million}. \tabularnewline
Formation & \en{Education} --- diplômes, dates, établissement. \tabularnewline
Compétences & \en{Skills} --- langues avec niveau CECRL (\en{French: native; English: B2}), informatique, permis. \tabularnewline
Références & \en{References available on request} --- ne listez pas de noms. \tabularnewline
Longueur & 1 à 2 pages maximum. \tabularnewline
\end{tabularx}
\end{tableau}

\subsection{Modèle 6 --- Essai argumentatif en 5 paragraphes}
\begin{dialoguebox}[Structure classique exigée dans les examens]
\begin{anglais}
\textbf{1. Introduction :} amorce + reformulation du sujet + annonce du plan.\\
\en{In recent years, \ldots has become a topic of debate. While some argue that \ldots,
others believe that \ldots\ This essay will examine both views before reaching a conclusion.}

\textbf{2. Argument 1 :} idée + explication + exemple.\\
\en{Firstly, \ldots\ This is because \ldots\ For example, \ldots}

\textbf{3. Argument 2 :} même structure.\\
\en{Secondly, \ldots\ In addition, \ldots\ A clear example of this is \ldots}

\textbf{4. Contre-argument et réfutation :}\\
\en{Admittedly, some people claim that \ldots\ However, this argument overlooks the fact
that \ldots}

\textbf{5. Conclusion :} synthèse + opinion + ouverture.\\
\en{In conclusion, although \ldots, it seems clear that \ldots\ Ultimately, the question is
not whether \ldots but how \ldots}
\end{anglais}
\end{dialoguebox}

\subsection{Modèle 7 --- Compte rendu de réunion}
\begin{dialoguebox}[Structure]
\begin{anglais}
\textbf{Meeting minutes --- 12 May 2026, 10:00, Head Office}\\
\textbf{Present:} \ldots \quad \textbf{Apologies:} \ldots \quad \textbf{Chair:} \ldots

\textbf{1. Review of Q1 results.} Sales increased by 12\,\% compared with the same period
last year. Costs also rose, so the margin is slightly lower. Action: M. Diop to send the
detailed breakdown by 19 May.

\textbf{2. New market.} The team agreed to postpone the launch until September. Action:
F. Ndiaye to prepare a revised budget.

\textbf{3. Any other business.} A follow-up meeting will be held next week.

\textbf{Next meeting:} 26 May 2026, 10:00.
\end{anglais}
\end{dialoguebox}
\textbf{Formules fixes :} \en{Action: [nom] to [verbe] by [date]} ;
\en{The team agreed to\ldots} ; \en{It was decided that\ldots} ;
\en{Apologies:} (absents excusés) ; \en{Any other business} (questions diverses).

\subsection{Modèle 8 --- Message court (WhatsApp, SMS professionnel)}
\begin{tableau}[Registre court mais correct]
\begin{tabularx}{\linewidth}{@{}p{6.4cm}X@{}}
\entete{Message} & \entete{Emploi} \tabularnewline
\hline
\en{Hi Ali, are you free at 3? Need to talk about the report.} & Court, neutre, efficace \tabularnewline
\en{Running 10 min late, sorry!} & Retard --- acceptable même au travail \tabularnewline
\en{Got it, thanks!} & Accusé de réception \tabularnewline
\en{Could you send me the file when you have a moment?} & Demande polie sans urgence \tabularnewline
\en{Thanks again for today --- much appreciated.} & Remerciement après une réunion \tabularnewline
\en{Just a quick reminder about tomorrow's meeting at 10.} & Rappel poli \tabularnewline
\en{No worries, we'll sort it out tomorrow.} & Dédramatiser \tabularnewline
\end{tabularx}
\end{tableau}

\section{Partie E --- Le journal écrit : 60 invites}
Écrire cinq à dix phrases par jour est l'exercice le plus rentable de tout le livre : il
révèle immédiatement quels mots et quelles structures vous manquent.

\begin{tableau}[60 invites de journal (deux par jour pendant un mois)]
\footnotesize
\begin{multicols}{2}
\begin{anglais}
\begin{enumerate}\setlength{\itemsep}{1pt}
\item Describe your morning in detail.
\item What made you smile today?
\item Describe your bedroom to a blind person.
\item Write about someone who helped you this week.
\item What are you looking forward to?
\item Describe the best meal you have ever eaten.
\item What would you do with a free week?
\item Write about a fear you would like to overcome.
\item Describe your favourite place in your city.
\item What did you learn this week?
\item Write a letter to yourself in five years.
\item Describe the weather and how it affects you.
\item What habit would you like to change?
\item Write about your first job.
\item Describe a stranger you remember.
\item What are you grateful for today?
\item Write about a mistake that taught you something.
\item Describe your daily route to work.
\item What does success mean to you?
\item Write about a song that moves you.
\item Describe your family in five sentences each.
\item What would you change about your country?
\item Write about a time you were very tired.
\item Describe something you own and treasure.
\item What do you want to be remembered for?
\item Write about an argument you had.
\item Describe a market or a busy street.
\item What is the hardest thing about learning English?
\item Write about a person you admire from afar.
\item Describe a journey that went wrong.
\item What do you do when you are angry?
\item Write about a tradition in your family.
\item Describe your phone and what it does for you.
\item What would you tell a younger version of yourself?
\item Write about a book or film that changed you.
\item Describe the sound of your neighbourhood.
\item What are you avoiding doing? Why?
\item Write about a kindness you received.
\item Describe your ideal working day.
\item What is overrated in your opinion?
\item Write about a time you were lost.
\item Describe your hands and what they have done.
\item What do you believe but cannot prove?
\item Write about a celebration you attended.
\item Describe your relationship with money.
\item What skill would you like to have?
\item Write about a teacher who mattered.
\item Describe something you built or fixed.
\item What makes you laugh?
\item Write about a promise you made.
\item Describe the view from your window right now.
\item What would you do if you had no fear?
\item Write about a conversation you regret.
\item Describe your favourite time of day.
\item What do you need to let go of?
\item Write about a place you will never return to.
\item Describe your English level honestly.
\item What are your three goals for next month?
\item Write about someone you miss.
\item Summarise your month in ten sentences.
\end{enumerate}
\end{anglais}
\end{multicols}
\end{tableau}

\begin{astuce}[La méthode des trois jets]
\textbf{Jet 1 (5 min) :} écrivez sans vous arrêter, sans dictionnaire, sans corriger.
\textbf{Jet 2 (5 min) :} relisez et corrigez la grammaire, les temps, les articles ---
c'est là que vous apprenez. \textbf{Jet 3 (3 min, deux fois par semaine) :} faites corriger
par un outil ou un partenaire, et notez \textbf{une seule} erreur récurrente à surveiller
la semaine suivante. Une erreur corrigée à la fois : c'est la seule méthode qui fonctionne.
\end{astuce}

\begin{exo}[17.1 --- Orthographe]
Écrivez la forme correcte :
\en{stop (past), plan (-ing), study (3rd person), happy (comparative), leaf (plural),
life (plural), travel (past, BrE), colour (AmE), centre (AmE), practise (noun, BrE),
organise (AmE), programme (AmE)}.
\end{exo}
\begin{corrige}[17.1]
\en{stopped, planning, studies, happier, leaves, lives, travelled, color, center, practice,
organize, program}.
\end{corrige}

\begin{exo}[17.2 --- Ponctuation]
Corrigez ces cinq phrases (ponctuation et majuscules) :
\begin{enumerate}
  \item \en{*hello ! how are you today ?}
  \item \en{*i was born in 12.03.1995 in thies}
  \item \en{*he said ``i'm tired'' , and then he left}
  \item \en{*the meeting is on monday at 14.30}
  \item \en{*i bought apples , bananas and oranges}
\end{enumerate}
\end{exo}
\begin{corrige}[17.2]
1. \en{Hello! How are you today?} 2. \en{I was born on 12 March 1995 in Thiès.}
3. \en{He said, ``I'm tired,'' and then he left.} 4. \en{The meeting is on Monday at 2.30 pm.}
5. \en{I bought apples, bananas, and oranges.}
\end{corrige}

\begin{exo}[17.3 --- Écrivez]
Rédigez (150 à 200 mots) : un e-mail formel à une école de langue pour demander les dates,
les tarifs et les modalités d'inscription. Utilisez au moins :
\en{I am writing to\ldots}, \en{I would be grateful if you could\ldots},
\en{Please find attached\ldots}, \en{I look forward to hearing from you},
\en{Yours faithfully}.
\end{exo}
\begin{corrige}[17.3 --- exemple de réponse attendue]
\begin{anglais}
\textbf{Subject: Enquiry about English courses --- September intake}

Dear Sir or Madam,

I am writing to enquire about the intensive English courses advertised on your website.

I would be grateful if you could let me know the start dates for the September intake, the
weekly timetable, and the total cost of the course. I would also like to know whether an
entrance test is required, and whether payment can be made in two instalments.

I have been studying English on my own for one year, and I believe my level is
approximately B1. Please find attached a short description of my goals, in case it is useful
for the placement test.

I look forward to hearing from you.

Yours faithfully,\\
Awa Sarr
\end{anglais}
\end{corrige}

\begin{mission}[Chapitre 17 --- lire et écrire chaque jour]
\textbf{Tous les jours :} 15 minutes de lecture extensive (facile, sans dictionnaire) +
10 minutes de journal (5 phrases minimum).\\
\textbf{Deux fois par semaine :} un texte de ce livre relu à voix haute, puis réécrit de
mémoire --- exercice puissant qui installe les structures.\\
\textbf{Une fois par semaine :} un e-mail complet (modèle 2, 3 ou 4) envoyé à un vrai
destinataire --- partenaire de langue, école, entreprise. L'écrit envoyé progresse trois
fois plus vite que l'écrit dans un cahier.
\end{mission}
